\section{A multilevel code for two-step reconstructions}
\label{app:mlmc-code}

This appendix gives another finite class of representatives for the
original robustness support. Its members are finite programs using
arithmetic, tanh, Euclidean projection, and a fixed finite median rule.
They need not themselves be robust networks. All randomness below is
used only to prove the existence of a representative in a fixed finite
class; the representative is deterministic after its coefficients and
finite computation graph have been recorded.

\begin{theorem}[Multilevel finite code]
\label{thm:mlmc-code}
There is a universal constant $A$ with the following property. Let
$d,H\ge1$, $K>0$, and $0<r\le1/6$, and put
\[
 E=K+H+1,\qquad
 T_0=d+\log(2H+2),\qquad
 \Lambda=\log(2+d+E+r^{-1}).
\]
There is a finite class $\mathcal G_{K,r,d,H}$, independent of learner
width, such that
\begin{equation}\label{eq:mlmc-budget}
 \log|\mathcal G_{K,r,d,H}|
 \le A T_0\left[\frac{E^3}{r^3}
                  +\frac{(d+1)E^2}{r^2}\right]\Lambda.
\end{equation}
For every width $N$ and every network $w$ with $\kappa_N(w)\le K$,
there is $g\in\mathcal G_{K,r,d,H}$ whose output is the output of the
original network with input-dependent disturbances of stacked norm at
most $r$. Consequently, if $w\in\mathcal R_{N,c}$ and
$r\le c/\sqrt{1+HE_v(w)}$, then
\[
 \sup_{x\in\{-1,1\}^d}|T(z_w(x))-T(g(x))|\le1/16.
\]
\end{theorem}

We prove the theorem in three steps. First we construct an economical
code for a map containing two activations. We then embed such maps in
a virtual computation of the original network. Finally we choose
spectral cutoffs and error allocations jointly across layers.

\subsection{A two-activation sampling lemma}

Consider the vector-valued map
\begin{equation}\label{eq:mlmc-map}
 F(z,z')=\sum_{i=1}^N w_i\tanh\left(a_i\cdot z+
                \sum_{j=1}^N R_{ij}\tanh(c_j\cdot z')\right),
 \qquad w_i\in\R^{q_o},
\end{equation}
where $a_i\in\R^{q_a}$ and $c_j\in\R^{q_c}$. Assume
\begin{equation}\label{eq:mlmc-map-energy}
 N\sum_i\norm{w_i}_2^2\le k_o,\quad
 \norm{R}_F^2\le k_R,\quad
 \frac1N\sum_i\norm{a_i}_2^2\le k_a,\quad
 \frac1N\sum_j\norm{c_j}_2^2\le k_c.
\end{equation}
Write
\[
 A_i=\sum_j|R_{ij}|,\qquad M_i=\max\{1,A_i\},\qquad
 B=\sum_i\norm{w_i}_2,
\]
\begin{equation}\label{eq:mlmc-work-resources}
 D=\sum_i\norm{w_i}_2M_i,\qquad
 Z=\sum_i\frac{\norm{w_i}_2}{M_i}.
\end{equation}
Cauchy--Schwarz gives
\begin{equation}\label{eq:mlmc-work-energy}
 B\le\sqrt{k_o},\qquad B^2\le ZD,\qquad
 D^2\le2k_o(1+k_R).
\end{equation}
Indeed, $\sum_i A_i^2\le N\norm{R}_F^2$, so
$D\le\sqrt{k_o}(1+\sqrt{k_R})$.

\begin{lemma}[Antithetic map sampling]
\label{lem:mlmc-map-sampling}
Fix $\tau>0$ and a finite collection of evaluation arguments. There
is a random finite program approximating \eqref{eq:mlmc-map} with the
following properties. At each fixed argument, its error is at most
$9\tau$ except on an event of probability at most $e^{-c_1g}$, where
$g$ is the odd number of independent median groups. Its expected number
of stored scalar coefficients is at most
\begin{equation}\label{eq:mlmc-map-work}
 c_2g(1+q_a+q_c+q_o)
                 \frac{k_o(1+k_R)}{\tau^2}.
\end{equation}
After an event of constant probability bounding total coefficient work,
the program can be rounded to a fixed grid with one further error at
most a constant times $\tau$. All parameter ranges and the inverse
grid mesh are polynomial in the energy and dimension upper bounds,
the evaluation-argument norm bounds, $\tau^{-1}$, and the work bound.
No original-width neuron index is needed in the stored program.

The probability assertion also holds conditionally when the evaluation
arguments have been produced by earlier independent random programs.
\end{lemma}

\begin{proof}
We give the sampling, coefficient, and rounding details. Deterministic
feature pruning will first ensure finite coefficient ranges. Removing
roots with $\norm{a_i}_2>B_a$ changes the map by at most
\[
 \sum_{\norm{a_i}_2>B_a}\norm{w_i}_2
       \le\sqrt{k_ok_a}/B_a.
\]
Deleting columns $j$ with $\norm{c_j}_2>B_c$ from $R$ changes it by at
most
\begin{align*}
 \sum_i\norm{w_i}_2\sum_{j:\norm{c_j}_2>B_c}|R_{ij}|
 &\le\sqrt{k_o/N}\,
       \sqrt{\#\{j:\norm{c_j}_2>B_c\}\,k_R}\\
 &\le\sqrt{k_ok_Rk_c}/B_c.
\end{align*}
Choose the two cutoffs to make these errors at most $\tau$ each;
zero numerators require no cutoff. All old energy upper bounds remain
valid after pruning. Recompute $A_i,M_i,B,D,Z$ for the pruned map.
If $B\le\tau$, the zero map already has the required accuracy, so
assume $B>\tau$.

For $A_i>0$, sample $J$ with probability $|R_{iJ}|/A_i$ and put
\[
 X_i=A_i\operatorname{sgn}(R_{iJ})\tanh(c_J\cdot z').
\]
Then $|X_i|\le A_i$ and
$\E X_i=\mu_i:=\sum_jR_{ij}\tanh(c_j\cdot z')$.
For $A_i=0$, use $X_i=0$. Let
\[
 n_{i,0}=\lceil M_i^2\rceil,\qquad n_{i,l}=2^l n_{i,0}.
\]
At level zero define
$\Delta_{i,0}=\tanh(a_i\cdot z+\overline X_{i,n_{i,0}})$.
At level $l\ge1$, split $n_{i,l}$ independent observations into two
equal halves with averages $Y,Y'$, and define
\begin{equation}\label{eq:mlmc-antithetic}
 \Delta_{i,l}=
 \tanh\left(a_i\cdot z+\frac{Y+Y'}2\right)
 -\frac12\tanh(a_i\cdot z+Y)
 -\frac12\tanh(a_i\cdot z+Y').
\end{equation}
All three evaluations reuse the same child activations. The bounded
second derivative of tanh and the fourth moment of a bounded sample
average imply, with universal constants,
\begin{equation}\label{eq:mlmc-level-estimates}
 \E|\Delta_{i,l}|^2\le c_3 2^{-2l},\qquad
 \left|\E\tanh(a_i\cdot z+\overline X_{i,n_{i,l}})
             -\tanh(a_i\cdot z+\mu_i)\right|
       \le c_4 2^{-l}.
\end{equation}
For completeness, Taylor's theorem bounds the absolute value in
\eqref{eq:mlmc-antithetic} by
$\norm{\tanh''}_\infty|Y-Y'|^2/8$.
The fourth moment of a centered average of $m$ independent variables
bounded by $2A_i$ is at most a universal constant times $A_i^4/m^2$.
Since $n_{i,0}\ge A_i^2$, this proves the first estimate. The second
follows from the same Taylor bound and the variance of the average.
The estimates are uniform in both evaluation arguments.

Choose a finite last level $l_*$ and let
\[
 p_i=\frac{\norm{w_i}_2}{M_iZ},\qquad
 \pi_l=\frac{2^{-3l/2}}{\sum_{j=0}^{l_*}2^{-3j/2}}.
\]
Zero-weight roots are omitted. Independently choose a root and level,
and return the vector $Y=w_i\Delta_{i,l}/(p_i\pi_l)$.
The antithetic expectations telescope. Therefore, if $F_*=\E Y$,
\begin{equation}\label{eq:mlmc-root-estimates}
 \norm{F_*-F}_2\le c_4B2^{-l_*},\qquad
 \E\norm{Y}_2^2\le c_5ZD,
 \qquad \E\operatorname{work}(Y)\le c_6D/Z.
\end{equation}
Here scalar work counts child activations, with a fixed constant for
each root. The identities proving the last two estimates are
\[
 \sum_i\frac{\norm{w_i}_2^2}{p_i}=ZD,\quad
 \sum_i p_iM_i^2=D/Z,\quad
 \sum_l2^{-2l}/\pi_l=O(1),\quad
 \sum_l\pi_l2^l=O(1).
\]
Choose $2^{l_*}\ge c_4\sqrt{k_o}/\tau$, or a larger common polynomial
upper bound, making the bias at most $\tau$.

Average $m=\lceil c_7ZD/\tau^2\rceil$ independent copies of $Y$ in
each group. For sufficiently large $c_7$, Chebyshev's inequality makes
the probability of vector error exceeding $\tau$ about $F_*$ at most
$1/64$. The expected work of a group is
\[
 O(D^2/\tau^2+D/Z)=O(D^2/\tau^2),
\]
because $B>\tau$ and $B^2\le ZD$ imply
$D/Z\le D^2/B^2<D^2/\tau^2$.

Use $g$ independent groups and the following fixed median rule: choose
the first group output with at least $(g+1)/2$ group outputs within
distance $4\tau$, and return zero if there is no such group. Except
with probability $e^{-c_1g}$, a strict majority of the groups lie
within $\tau$ of $F_*$. Every such group qualifies, and every selected
group lies within $5\tau$ of $F_*$. If each group output is subsequently
changed by at most $\tau$, the same rule returns a point within
$6\tau$ of $F_*$. This follows because the qualifying ball must
intersect the strict majority of good groups. Thus no Lipschitz property
of the median operation is needed. Combining the two pruning errors,
finite-level bias, and the rounded median gives error at most $9\tau$.

Multiplying scalar work by $1+q_a+q_c+q_o$ bounds the number of recorded
scalar coefficients: record each selected root feature, each child
feature, its signed inner coefficient, and the vector coefficient of
each term in a group average. This proves \eqref{eq:mlmc-map-work}.
For several maps, apply Markov's inequality to their \emph{total
dimension-weighted coefficient work}, rather than separately to their
unweighted child counts. On an event of constant probability, this
total is at most its chosen common deterministic bound $W$. Every
selected root then has $M_i^2\le n_{i,l}\le W$.

The vector coefficient of an averaged term is bounded by
\[
 \norm{\frac{w_i}{mp_i\pi_l}}_2
 \le\frac{\tau^2M_i}{c_7D\pi_l}
 \le\frac{\tau\sqrt W}{c_7\pi_l},
\]
using $D\ge B>\tau$. Inner coefficients have magnitude at most
$A_i\le\sqrt W$. Feature pruning bounded the remaining feature
coefficients, and the finite-level inverse probabilities are polynomial
in the common resources. Thus every recorded coefficient lies in a
fixed polynomially bounded interval.

Suppose all evaluation arguments have norm at most a fixed polynomial
bound. There are at most $W$ stored coefficients and only two nested
activation levels. The 1-Lipschitz property of tanh shows that rounding
all coefficients to a sufficiently fine reciprocal-polynomial mesh
changes every group output by at most $\tau$, uniformly on these
arguments. The preceding median argument applies after this rounding.
The forest topology, its group assignment, and its half-list divisions
have description length $O(W\log(2+W))$. Only numeric feature vectors
are stored, so no original-width neuron label appears. Finally, all
stochastic estimates were uniform at a fixed argument. Conditioning on
arguments produced by independent preceding programs therefore preserves
them, and a union bound applies to any finite collection of arguments.
\end{proof}

\subsection{Realization inside the original robustness ball}

Write $|h|_N=\norm{h}_2/\sqrt N$ and let $k_s=\norm{W^s}_F^2$ for
$2\le s\le H$, with $k_{H+1}=E_v(w)$. Split each internal matrix by
singular value truncation at a threshold $a_s\in(0,1)$:
\[
 W^s=A_s+R_s,\qquad A_s=U_s\Sigma_sV_s^\top,\qquad
 \norm{R_s}_{\rm op}\le a_s,\qquad
 q_s=\operatorname{rank}A_s\le k_s/a_s^2.
\]
For drive coordinates $t_s\in\R^{q_s}$ use the original-space vector
$u_s=\sqrt N U_st_s$. An exact low-rank drive of a tanh state $h$ is
$t_s=\Sigma_sV_s^\top h/\sqrt N$, of norm at most $\sqrt{k_s}$.

The first drive is the one-activation map
$\Sigma_2V_2^\top h^1(x)/\sqrt N$. The next drive uses $h^1(x)$ as its
left boundary. All later drive maps have the two-activation form
\begin{equation}\label{eq:mlmc-drive-map}
 \frac1{\sqrt N}\Sigma_sV_s^\top
 \tanh\left(\sqrt N U_{s-1}t_{s-1}
       +R_{s-1}\tanh(\sqrt N U_{s-2}t_{s-2})\right).
\end{equation}
In Lemma~\ref{lem:mlmc-map-sampling}, the output-energy bound is $k_s$,
the residual-energy bound is $k_{s-1}$, and the root and child feature
second moments are $q_{s-1}$ and $q_{s-2}$. For an early map using the
original first state, its child features are rows of $W^1$, whose
normalized squared norm sum is at most $K/d$. The final readout has
the same form with scalar output and output-energy bound $k_{H+1}$.

Construct and round each map while advancing the computation. After
each drive map, project its output onto the ball of radius
$\sqrt{\max(K,H,1)}$. This does not increase error about its exact
target drive. Let $\widehat t_s$ be the resulting coordinates and
define full virtual states by
\[
 \widehat h^1=h^1,\qquad
 \widehat h^2=\tanh(\sqrt N U_2\widehat t_2+R_2h^1),
\]
\begin{equation}\label{eq:mlmc-virtual-state}
 \widehat h^s=\tanh\left(\sqrt N U_s\widehat t_s+
                  R_s\tanh(\sqrt N U_{s-1}\widehat t_{s-1})\right),
 \qquad s\ge3.
\end{equation}
These full states are used only to verify the disturbance bound; the
stored program advances the drive coordinates through its finite map
codes and evaluates the final readout code.

Suppose the drive-map errors are at most $\delta_s$, and the final
readout-map error is at most $\delta_{H+1}$. Their exact targets are
the low-rank drives of the already rounded preceding virtual states.
For $s\ge3$, comparison with the original recurrence gives
\[
 |\widehat h^s-\tanh(W^s\widehat h^{s-1})|_N
 \le\delta_s+a_sa_{s-1}.
\]
Indeed, the low-rank preactivation discrepancy has normalized norm
at most $\delta_s$, while the discarded older residual passes through
$R_{s-1}$ and $R_s$, with a 1-Lipschitz tanh between. Its left boundary
is a tanh state of normalized norm at most one. At stage two the bound
is just $\delta_2$; the first-state disturbance is zero. Consequently
the original-network disturbances realizing the virtual output satisfy
\begin{equation}\label{eq:mlmc-disturbance}
 \norm\eta_*
 \le\left(\sum_{s=2}^{H+1}\delta_s^2\right)^{1/2}
   +\left(\sum_{s=3}^H a_s^2a_{s-1}^2\right)^{1/2}.
\end{equation}
For $H=1$, only the one-activation readout map is needed; for $H=2$,
the exact first-state boundary is retained throughout.

Every stage uses fresh map samples independent of earlier map programs.
Lemma~\ref{lem:mlmc-map-sampling} therefore applies conditionally at
the already rounded prefix. Choose $g=O(T_0)$ so a union bound over
all $2^d$ inputs and all maps fails with probability at most $1/100$.
The coefficient-work event can be chosen to fail with probability at
most $1/100$, so the two events intersect. All rounding is already
included in the map errors in \eqref{eq:mlmc-disturbance}.

\subsection{Adaptive cutoffs and a common finite class}

Put $e_s=1+k_s$ for the internal blocks and the readout. Their sum is
at most $E=K+H+1$. For a sufficiently small universal constant $c_0$,
choose
\begin{equation}\label{eq:mlmc-cutoffs}
 a_s^2=\frac{c_0r\sqrt{e_s}}{\sqrt E},\qquad 2\le s\le H.
\end{equation}
Then $a_s<1$ and
\[
 \sum_{s=3}^H a_s^2a_{s-1}^2
 =\frac{c_0^2r^2}{E}\sum_{s=3}^H\sqrt{e_se_{s-1}}
 \le c_0^2r^2.
\]
With $\alpha=\sqrt E/(c_0r)$, the retained ranks satisfy
$q_s\le\alpha\sqrt{e_s}$.

Ignoring fixed constants, define each map's coefficient resource as
\begin{equation}\label{eq:mlmc-profile-resource}
 B_s=(1+q_{s-2}+q_{s-1}+q_s)k_s(1+k_{s-1}).
\end{equation}
At the first two maps, replace the missing input dimension by $d+1$
and omit a nonexistent residual factor. At the readout, replace the
output dimension by one and the output energy by $k_{H+1}$. These
conventions give exactly the resources in
Lemma~\ref{lem:mlmc-map-sampling}. Apart from the first two original-
input contributions, they are bounded by a universal constant times
\[
 \alpha e_se_{s-1}
    (\sqrt{e_{s-2}}+\sqrt{e_{s-1}}+\sqrt{e_s}).
\]
At boundary indices in this estimate, set a missing $e_j$ equal to one;
the finitely many auxiliary indices change the following bounds only
by universal constants.
For each of the three indicated shifts $j$,
\[
 \sum_s e_s^{1/2}e_{s-1}^{1/2}e_j^{1/4}
 \le(\max_j e_j)^{1/4}\sum_s\sqrt{e_se_{s-1}}
 \le c_8' E^{5/4}.
\]
Using $\sqrt{u+v+w}\le\sqrt u+\sqrt v+\sqrt w$, and treating the
two original-input terms separately, gives
\begin{equation}\label{eq:mlmc-profile-sum}
 \left(\sum_s\sqrt{\max(1,B_s)}\right)^2
 \le c_8\left[\frac{E^3}{r}+(d+1)E^2\right].
\end{equation}
Inserting the ones increases the unsquared sum by at most $H+1$,
so its additional squared contribution is $O(H^2)$, already absorbed
by the right side. Each original-input term contributes at most a constant
times $\sqrt{d+1}\,E$ before squaring.

Let $S_0$ be the square root of the right side of
\eqref{eq:mlmc-profile-sum}, enlarged by a fixed constant. Set
\[
 \delta_s^2=\frac{r^2}{64S_0}\sqrt{\max(1,B_s)},
\]
and round each $\delta_s$ downward to a power-of-two multiple of a fixed
reciprocal-polynomial base, losing at most a factor two. Thus
$\sum_s\delta_s^2\le r^2/64$, every nonzero tolerance has a common
reciprocal-polynomial lower bound, and
\begin{equation}\label{eq:mlmc-total-work}
 \sum_s B_s\delta_s^{-2}
 \le c_9\left[\frac{E^3}{r^3}
                    +\frac{(d+1)E^2}{r^2}\right].
\end{equation}
Apply the sampling lemma with $\tau_s$ a sufficiently small fixed
multiple of $\delta_s$. Equations \eqref{eq:mlmc-disturbance} and
\eqref{eq:mlmc-cutoffs} put the complete virtual computation inside
the radius-$r$ disturbance ball.

To make the code common to all widths and all retained networks, record
the integer rank profile, the dyadic tolerance choices, and the finite
map forests. Every rank is at most $E/(c_0r)$. The bounds
\eqref{eq:mlmc-total-work} and $g=O(T_0)$ give a common polynomial bound
$W$ on total coefficient work on the chosen Markov event. Feature
cutoffs, coefficient magnitudes, inverse grid mesh, and finite last
levels are all polynomial in $2+d+E+r^{-1}$, with universal degrees:
this follows from the common rank bound, the common lower tolerance
bound, and the coefficient estimates in the sampling lemma. Fix those
ranges and that grid before defining the class of all recorded programs.
The possible topology, ranks, and tolerance choices cost
$O(W\Lambda+H\Lambda)$ nats; the coefficient grid costs $O(W\Lambda)$.
Since $H\le E$ and $r\le1/6$, the former additional term is absorbed.
This proves \eqref{eq:mlmc-budget}. The class contains no learner-width labels,
so the same class works for every $N$. This completes the proof of
Theorem~\ref{thm:mlmc-code}.

\subsection{The original posterior at the stronger scales}
\label{sec:mlmc-posterior}

The new code changes neither the prior nor the conditioning event.
The following statement records the exact scale hypotheses used with it.
Use the product-cylinder cost
$D_*=D_*^{\rm prod}$ from \eqref{eq:product-cylinder-cost}, and
write
\[
 P_N=N(d+1)+(L-2)N^2+N.
\]
For the chosen energy cutoff $K$, retain the radius
\[
 r=\min\{1/6,c/\sqrt{1+HK}\}.
\]
Let $J_{\rm ML}$ be the upper bound in \eqref{eq:mlmc-budget}, and let
$J$ be the earlier dense-code budget in \eqref{eq:budget}, with this same
$K$ and $r$. Put
\[
 J_{\rm best}=\min\{J,J_{\rm ML}\}.
\]
This minimum is legitimate without a union penalty: choose, from the
known resource bounds, whichever of the two common finite classes has
the smaller stated cardinality bound.

\begin{theorem}[A mixed width--depth bound at the original stronger scales]
\label{thm:mlmc-posterior}
Fix $0<\varepsilon,\delta<1$. Suppose the target has the Boolean
implementation specified in the paper. The sample condition
\[
 n\ge\frac2\varepsilon
       \left[J_{\rm best}+\log(1/\delta)\right]
\]
implies, with sample probability at least $1-\delta$,
$R(Q_{S,N,t,c})\le\varepsilon$ simultaneously at every indicated width
and scale in either of the following regimes:
\begin{enumerate}
\item Set $K=8D_*$. Allow every $N\ge C$ and every positive scale with
\[
 t^2\le\min\left\{1,
 \frac{ND_*}{(P_N/2)\log2+\log(2/\varepsilon)}\right\}.
\]
\item Set
\[
 K=8D_*+4(d+L)+4\log(2/\varepsilon).
\]
Allow every $N\ge C$ and $0<t\le N^{-1/2}$, in particular every
power-law scale with $\gamma\ge2$.
\end{enumerate}
In the first regime, at each fixed width the same sample event gives
$\limsup_{t\downarrow0}R(Q_{S,N,t,c})\le\varepsilon$.
At fixed $d,c,\varepsilon,\delta$, a sufficient sample count is
\[
 \widetilde O\!\left(
      \min\{C^4L^7,C^{9/2}L^6\}\right).
\]
\end{theorem}

\begin{proof}
Every robust network with $\kappa_N(w)\le K$ has calibrated radius at
least $r$. Theorem~\ref{thm:mlmc-code}, or the earlier dense code,
therefore supplies a representative with clipped-score error at most
$1/16$ at every Boolean input. The representative need not be a robust
network; only its membership in a common finite class is used next.

Use the ramp loss that is one at signed clipped scores at most $1/3$
and zero at scores at least $5/12$. Every representative of a fitting
network has zero empirical ramp loss, because its signed clipped
training score is at least $7/16$. At every test-margin failure of the
original network, its representative has signed clipped score at most
$5/16$, so the ramp loss is one. A union bound over the chosen finite
class shows that, on a single sample event of probability at least
$1-\delta$, every retained fitting network has test-margin error at
most $\varepsilon/2$. This sample event is independent of learner width
and Gaussian scale.

The product-cylinder lower bound and the Gaussian exponential moment give
the same tail as \eqref{eq:tail}, with the current valid cylinder cost:
\[
 Q_{S,N,t,c}\{\kappa_N>K\}
 \le\exp\!\left[-\frac{NK}{4t^2}
       +\frac{P_N}{2}\log2+\frac{ND_*}{t^2}\right].
\]
In the first regime the displayed scale inequality makes this at most
$\varepsilon/2$. In the second regime use
$t^2P_N/N\le d+L$ and the stated value of $K$, exactly as in the proof
of Theorem~\ref{thm:main}. Adding the discarded mass proves the risk
claim. At fixed $N$, the first scale inequality holds for all sufficiently
small positive $t$, proving the stated strong-prior conclusion without
identifying a separately defined zero-temperature posterior.

Finally, the product-cylinder calculation gives
$K=\widetilde O(CL)$ and $r^{-2}=\widetilde O(CL^2)$.
The budget \eqref{eq:mlmc-budget}, with $K+H+1=\widetilde O(CL)$,
is therefore $\widetilde O(C^{9/2}L^6)$ at the fixed resources stated
in the theorem. Taking the better of this and the existing
$\widetilde O(C^4L^7)$ code gives the final display.
\end{proof}

\begin{corollary}[Explicit logarithms in the mixed-power budget]
\label{cor:mlmc-powers}
Fix $d,c,\varepsilon,\delta$, and put
$F_0=1+\log^2(2+C)$ and $G=\log(2+CL)$.
In either regime of Theorem~\ref{thm:mlmc-posterior}, a sufficient
sample count is
\[
 O\!\left(G\min\left\{
 C^4L^7F_0^4\log^2(2+L),\;
 C^{9/2}L^6F_0^{9/2}\log(2+L)
 \right\}\right).
\]
In particular, at fixed teacher width the new branch is
$O(L^6\log^2(2+L))$.
\end{corollary}

\begin{proof}
The product-cylinder cost satisfies $D_*^{\rm prod}=O(CLF_0)$.
For either stated energy cutoff,
\[
 E=O(CLF_0),\qquad r^{-2}=O(CL^2F_0),\qquad
 T_0=O(\log(2+L)),\qquad \Lambda=O(G).
\]
Consequently the leading term in \eqref{eq:mlmc-budget} is
\[
 O\!\left(C^{9/2}L^6F_0^{9/2}\log(2+L)G\right).
\]
The term $(d+1)E^2r^{-2}$ has order $O(C^3L^4F_0^3)$ and is smaller
at the fixed ancillary resources. The dense-code calculation with the
same product-cylinder substitution gives
$O(C^4L^7F_0^4\log^2(2+L)G)$, as recorded after
Corollary~\ref{cor:all-scales-product-cylinder}. Taking the better code
and applying Theorem~\ref{thm:mlmc-posterior} proves the result.
\end{proof}

The second branch improves the stated depth dependence at fixed teacher
width, and is smaller at the level of the displayed powers when
$C<L^2$. It does not establish a better uniform quartic-width exponent.
No bounded spectral product, column coherence, additional output-margin
condition, or modified robustness support is used.
